\chapter{Introduction}
\label{ch:Introduction}
\pagenumbering{arabic}
The introduction shall give the reader an overview of the topic. 

This document is intended to serve as a template for a master's thesis written at the Institute for Mechanics and Fluiddynamics of the TU Bergakademie Freiberg. It also provides useful examples of how to use tables, figures, citations, and cross-references. Furthermore, it offers many tips and hints on how to write a proper thesis, regarding time management, the use of software, organization of data, code, etc. 

\section{Problem statement}
\label{sec:ProblemStatement}
Here, the specific topic of the thesis should be explained. Use the task sheet provided by your supervisors.

Students often struggle to write a proper thesis. The main difficulties are:
\begin{itemize}
    \item Structure and organization of the document
    \item Inconsistent mathematical notation
    \item Bad graphical representation of illustrations, diagrams and flow charts
    \item Incorrect referencing of chapters, sections, illustrations
    \item Citation of references (BibTex)
    \item Organization of data, code, and scripts
\end{itemize}
Another difficulty is the general format of the thesis. Several inconsistencies appear between templates provided by TUBAF and templates that were used previously for the IMFD theses.

\section{Things to do}
\begin{itemize}
    \item Title page
    \item Cover page (Berichte des Institutes für Mechanik und Fluiddynamik)
\end{itemize}

\section{Notation}
It is important to have a consistent mathematical notation. In mechanics, we have to deal with a lot of different quantities. Especially when different theories are used for which competing nomenclatures exist, it is challenging to develop a consistent nomenclature.

\subsection{Symbolic notation using different fonts}
Symbolic notation is very compact, which makes it easy for the reader to understand physical or mathematical concepts. The symbolic notation is independent on a specific reference system, which makes it very general. On the other hand, only a limited number of different fonts are available to distinguish between scalars, vectors, tensors of different orders, matrices, sets, configurations, and other classes that might be needed. Especially for PhD thesis, this often leads to difficulties in finding a consistent notation, due to the large number of symbols needed for different theories.
\begin{itemize}
    \item Scalars ($a$, $b$, \dots) are small letters
    \item Descriptive indices, operators ($\mathrm{a}$, $\mathrm{b}$, \dots) are small mathrm letters
    \item Scalars ($\alpha$, $\beta$, $\gamma$, $\lambda$, $\sigma$, $\varepsilon$) small greek letters
    \item Vectors ($\boldsymbol{x}$, $\boldsymbol{u}$$, \boldsymbol{\alpha}$, $\boldsymbol{\beta}$) bold small letters
    \item Vectors ($\boldsymbol{\mathrm{x}}$, $\boldsymbol{\mathrm{u}}$) small bold mathrm letters (only Latin alphabet)
    \item Second order tensors ($\boldsymbol{\sigma}$, $\boldsymbol{\varepsilon}$, $\boldsymbol{F}$) large bold letters (stress and strain are exceptions)
    \item Second orders tensors ($\boldsymbol{A}$, $\boldsymbol{B}$, $\boldsymbol{C}$, ) large bold letters
    \item Second orders tensors ($\boldsymbol{\mathrm{A}}$, $\boldsymbol{\mathrm{B}}$, $\boldsymbol{\mathrm{C}}$, ) large bold mathrm letters (only Latin alphabet)
    \item Third order tensor ($\mathfrak{a}$, $\mathfrak{E}$) for micromorphic theories 
    \item Third order tensor ($\boldsymbol{\mathfrak{a}}$, $\boldsymbol{\mathfrak{E}}$) bold for micromorphic theories 
    \item Forth order tensors ($\mathbb{C}$, $\mathbb{D}$, $\mathbb{Z}$, $\mathbb{R}$, $\Bbbk$)
    \item Sets ($\mathcal{A}$, $\mathcal{B}$, $\mathcal{F}$)
    \item Sets ($\boldsymbol{\mathcal{A}}$, $\boldsymbol{\mathcal{B}}$, $\boldsymbol{\mathcal{F}}$) bold
    \item Matrices ($\mathsf{A}$, $\mathsf{B}$) (assembly, general matrix notation)
    \item Matrices ($\boldsymbol{\mathsf{A}}$, $\boldsymbol{\mathsf{B}}$) bold (assembly, general matrix notation)
    \item Vectors ($\mathsf{f}$, $\mathsf{u}$) (force vector, displacement vector in general matrix notation)
    \item Vectors ($\boldsymbol{\mathsf{f}}$, $\boldsymbol{\mathsf{u}}$) bold (force vector, displacement vector in general matrix notation)
    \item Configurations ($\boldsymbol{X}$, $\boldsymbol{x}$) (initial, intermediate, deformed)
    \item Micro, Macro ($\boldsymbol{\sigma}$, $\bar{\boldsymbol{\sigma}}$ or $\boldsymbol{\sigma}$, $\boldsymbol{\Sigma}$,)
    \item mathtt ($\mathtt{a}$, $\mathtt{A}$)
    \item What about standard symbols for physical quantities (\href{https://en.wikipedia.org/wiki/List_of_physical_quantities}{Wiki})
\end{itemize}

Operators:
\begin{itemize}
    \item time derivative $\dot{\bullet}=\frac{\bullet}{\textrm{d}t}$
    \item volume average $\bar{\bullet}=\frac{1}{V} \int_V \bullet \, \textrm{d}V$
    \item \dots
\end{itemize}

\subsection{Symbolic notation using over- and underline accents}
\emph{Symbolic notation} can be extended using over- and underline accents to distinguish between, scalars, vectors, and tenors. Or between local and global quantities. Furthermore, descriptive indices can be used to define quantities belonging to different configurations or classes. 


\subsection{Index notation}

\section{Structure of the thesis}
Explain the structure of the thesis (approx. 1 page).

The subsequent chapters of this thesis are organized as follows: Chapter \ref{ch:Examples} gives a comprehensive overview about the theories used in this thesis.
The main chapter explains the methods used to write a proper thesis. It focuses on the specific problems outlined in Sec.~\ref{sec:ProblemStatement}.
%
